\documentclass[10pt,a4paper]{article}

%double si on veut une version prof et une version élève. simple si on veut une seule version.
\newcommand{\buildmode}{simple}

\InputIfFileExists{version.tex}{}{}
% Version (définie par le script)
\providecommand{\version}{eleve}

\usepackage{feuilleinterro}

%%%à modifier à chaque fois

\renewcommand{\niveau}{2nde}
\renewcommand{\chapitre}{Équations, solution d'inéquations}
\renewcommand{\numchap}{0.3}
\renewcommand{\theme}{Nombres et calcul}

%%%titre "CORRECTION" au lieu de "INTERROGATION", sans les champs Nom/Prénom
\renewcommand{\interroversion}[1]{
    \noindent\textbf{\niveau}

    \vspace{-0.5cm}
    \begin{center}
        \textbf{\large CORRECTION -- \chapitre}
    \end{center}

    \vspace{0.1cm}
}

% Axe gradué en quarts, entre -2 et 2, où seuls 0 et 1 sont indiqués,
% avec les nombres à placer marqués en rouge.
% #1 : liste "valeur/étiquette" séparés par des virgules, ex. {-1/-1,0.5/\frac{1}{2}}
\newcommand{\axegraduepartiellesolution}[1]{
\begin{center}
\begin{tikzpicture}[scale=1.5]
\draw[->] (-2.5,0)--(2.5,0);
\foreach \x in {-2,..., 2}{
    \draw (\x,0.1)--(\x,-0.1);
}
\foreach \x in {-2,..., 1}{
    \draw (\x+0.5,0.07)--(\x+0.5,-0.07);
    \foreach \minx in {0.25, 0.5, 0.75}{
        \draw (\x+\minx, 0.05)--(\x+\minx, -0.05);
    }
}
\node[below] at (0,-0.1) {\(0\)};
\node[below] at (1,-0.1) {\(1\)};
\foreach \v/\lbl in {#1}{
    \filldraw[red] (\v,0) circle (1.5pt) node[above]{\scriptsize \(\lbl\)};
}
\end{tikzpicture}
\end{center}
}

\begin{document}

% =========================================================
% PAGE 1 : VERSION A
% =========================================================

\interroversion{A}

\textbf{Exercice 1 -- Nombres sur une droite graduée \hfill /3}

\axegraduepartiellesolution{-1/-1, 0.5/\frac{1}{2}, -1.75/-\frac{7}{4}}

\textbf{Exercice 2 -- Solution d'une équation \hfill /4}

Testons \(x=3\) dans \((E_1)\) :
\[
2\times3+1=6+1=7.
\]
On retrouve bien le membre de droite, donc \(\textcolor{blue}{3 \text{ est solution de } (E_1)}\).

Testons \(x=3\) dans \((E_2)\) :
\[
3\times3-2=9-2=7\neq10.
\]
Donc \(\textcolor{blue}{3 \text{ n'est pas solution de } (E_2)}\).

\textbf{Exercice 3 -- Résolution d'équation \hfill /3}

\[
\begin{aligned}
2x+1&=4\\
2x&=4-1\\
2x&=3\\
x&=\dfrac{3}{2}.
\end{aligned}
\]

L'ensemble des solutions est \(\textcolor{blue}{S=\left\{\dfrac{3}{2}\right\}}\).

\textbf{Exercice 4 -- Intervalles (bonus) \hfill /1}

La borne \(-3\) est incluse (\(\leqslant\)) et la borne \(2\) est exclue (\(<\)), donc l'intervalle recherché est
\[
\textcolor{blue}{\left[-3\ ;\ 2\right[}.
\]

\newpage

% =========================================================
% PAGE 2 : VERSION B
% =========================================================

\interroversion{B}

\textbf{Exercice 1 -- Nombres sur une droite graduée \hfill /3}

\axegraduepartiellesolution{2/2, -0.5/-\frac{1}{2}, 1.25/\frac{5}{4}}

\textbf{Exercice 2 -- Solution d'une équation \hfill /4}

Testons \(x=4\) dans \((E_1)\) :
\[
3\times4-5=12-5=7.
\]
On retrouve bien le membre de droite, donc \(\textcolor{blue}{4 \text{ est solution de } (E_1)}\).

Testons \(x=4\) dans \((E_2)\) :
\[
2\times4+3=8+3=11\neq13.
\]
Donc \(\textcolor{blue}{4 \text{ n'est pas solution de } (E_2)}\).

\textbf{Exercice 3 -- Résolution d'équation \hfill /3}

\[
\begin{aligned}
4x-1&=6\\
4x&=6+1\\
4x&=7\\
x&=\dfrac{7}{4}.
\end{aligned}
\]

L'ensemble des solutions est \(\textcolor{blue}{S=\left\{\dfrac{7}{4}\right\}}\).

\textbf{Exercice 4 -- Intervalles (bonus) \hfill /1}

La borne \(-1\) est exclue (\(<\)) et la borne \(4\) est incluse (\(\leqslant\)), donc l'intervalle recherché est
\[
\textcolor{blue}{\left]-1\ ;\ 4\right]}.
\]

\newpage

% =========================================================
% PAGE 3 : VERSION C
% =========================================================

\interroversion{C}

\textbf{Exercice 1 -- Nombres sur une droite graduée \hfill /3}

\axegraduepartiellesolution{-2/-2, -1.5/-\frac{3}{2}, 0.75/\frac{3}{4}}

\textbf{Exercice 2 -- Solution d'une équation \hfill /4}

Testons \(x=2\) dans \((E_1)\) :
\[
5\times2-3=10-3=7.
\]
On retrouve bien le membre de droite, donc \(\textcolor{blue}{2 \text{ est solution de } (E_1)}\).

Testons \(x=2\) dans \((E_2)\) :
\[
4\times2+1=8+1=9\neq17.
\]
Donc \(\textcolor{blue}{2 \text{ n'est pas solution de } (E_2)}\).

\textbf{Exercice 3 -- Résolution d'équation \hfill /3}

\[
\begin{aligned}
2x+5&=2\\
2x&=2-5\\
2x&=-3\\
x&=-\dfrac{3}{2}.
\end{aligned}
\]

L'ensemble des solutions est \(\textcolor{blue}{S=\left\{-\dfrac{3}{2}\right\}}\).

\textbf{Exercice 4 -- Intervalles (bonus) \hfill /1}

La borne \(-4\) est incluse (\(\leqslant\)) et la borne \(0\) est exclue (\(<\)), donc l'intervalle recherché est
\[
\textcolor{blue}{\left[-4\ ;\ 0\right[}.
\]

\newpage

% =========================================================
% PAGE 4 : VERSION D
% =========================================================

\interroversion{D}

\textbf{Exercice 1 -- Nombres sur une droite graduée \hfill /3}

\axegraduepartiellesolution{2/2, 1.5/\frac{3}{2}, -0.25/-\frac{1}{4}}

\textbf{Exercice 2 -- Solution d'une équation \hfill /4}

Testons \(x=5\) dans \((E_1)\) :
\[
2\times5-3=10-3=7.
\]
On retrouve bien le membre de droite, donc \(\textcolor{blue}{5 \text{ est solution de } (E_1)}\).

Testons \(x=5\) dans \((E_2)\) :
\[
3\times5+2=15+2=17\neq20.
\]
Donc \(\textcolor{blue}{5 \text{ n'est pas solution de } (E_2)}\).

\textbf{Exercice 3 -- Résolution d'équation \hfill /3}

\[
\begin{aligned}
4x+3&=-2\\
4x&=-2-3\\
4x&=-5\\
x&=-\dfrac{5}{4}.
\end{aligned}
\]

L'ensemble des solutions est \(\textcolor{blue}{S=\left\{-\dfrac{5}{4}\right\}}\).

\textbf{Exercice 4 -- Intervalles (bonus) \hfill /1}

La borne \(-2\) est exclue (\(<\)) et la borne \(3\) est incluse (\(\leqslant\)), donc l'intervalle recherché est
\[
\textcolor{blue}{\left]-2\ ;\ 3\right]}.
\]

\end{document}
